\section{The claimed construction, in outline}\label{sec:candidate}

This section follows the opening summary of the reconstructed source manuscript (file \fn{03\_}) and the monograph's orientation. Nothing in it was checked here, apart from the arithmetic noted.

\begin{enumerate}[label=\arabic*.]
\item \emph{Base.} The triangle group $\Delta=\Delta(3,4,\infty)\cong(\ZZ/3)*(\ZZ/4)$ acts on the upper half-plane. The quotient, with its cusp added, is $\Pp^1$ with orbifold points of orders $3$ and $4$ and one cusp.
\item \emph{A lattice representation.} On $V\cong\ZZ^4$ there are automorphisms $T_1$ of order $3$ and $T_2$ of order $4$ whose product has unipotent inverse $T_0$ with $(T_0-I)^2=0$; $\Delta$ acts on the dual lattice $\Lambda$ through them.
\item \emph{Periods and the family.} Period functions $\tau,\mu,\beta$ on the half-plane give an equivariant period map. The monograph writes the period matrix as
\[
\Pi=[Z\mid I_2],\qquad Z=\begin{pmatrix}6\mu&\tau\\ \beta&\mu\end{pmatrix},\qquad D=\operatorname{Im}\beta-\frac{6(\operatorname{Im}\mu)^2}{\operatorname{Im}\tau}<0,\qquad \det\nolimits_{\R}\Pi=\operatorname{Im}\tau\cdot D<0 .
\]
The quotient is a family of complex two-dimensional tori over the punctured orbifold curve. Its very general fibre is not algebraic, and its Néron--Severi group is generated by one indefinite class.
\item \emph{Three fillings.} At the two elliptic points, Kodaira's logarithmic transform: (disc $\times$ torus)$/(\ZZ/m)$, $m=3,4$, giving multiple fibres. At the cusp, Mumford's toric degeneration, whose central fibre is a degree-six del Pezzo surface glued to itself along opposite sides, with two triple points.
\item \emph{Gluing and the fundamental group.} The fillings are glued along collars, and the gluings can be changed by translations. The fundamental group and the homology depend on three integers $\ell_0,\ell_1,\ell_2$, with
\[
\pi_1(X)\cong\ZZ/|12\ell_0-4\ell_1-3\ell_2| .
\]
The record's candidate takes $(\ell_0,\ell_1,\ell_2)=(0,1,-1)$, for which $12\ell_0-4\ell_1-3\ell_2=-1$, so the fundamental group is trivial (source manuscript §7; monograph Part~I).
\item \emph{Conclusions claimed.} For this choice, the manuscript concludes:
\begin{itemize}
\item $X$ is simply connected with the integral homology of $S^6$;
\item $X$ is therefore diffeomorphic to $S^6$, as there are no exotic smooth structures on the six-sphere;
\item $c_3(X)=2$, localised at the cusp;
\item $a(X)=1$, with the fibration as algebraic reduction.
\end{itemize}
\end{enumerate}

The period data of item~3 are also what the owner's Yang--Mills and Erdős--Straus workbenches take from this construction (Section~\ref{sec:overlaps}). In the Yang--Mills record the quantity $D_{\mathrm{YM}}=L_Tq_T+6m_T^2$, with $L_T=\operatorname{Im}\tau$, $q_T=-\operatorname{Im}\beta$ and $m_T=\operatorname{Im}\mu$, equals $-\operatorname{Im}\tau\cdot D=-\det_{\R}\Pi>0$. So the two records use the same number, with opposite sign conventions.
